\section{Differentially stable polynomial ideals and normal cylinders}

All ideals in this chapter are the entire original polynomial ideals. Let
$R$ be a commutative $\mathbb Q$-algebra. The only scalar cancellation used
below is multiplication by the image of a nonzero rational inverse. No
reducedness or domain hypothesis on $R$ is needed.

\begin{lemma}[Every coefficient of a differential ideal remains in the ideal]
\label{hir:differential-stable-polynomial-coefficients}
\lean{Hironaka.HironakaDifferentialStableCylindricalIdeal.polynomial_C_coeff_mem_of_derivative_stable}
\leanok
Let $J\subseteq R[T]$ be an ideal such that $f'\in J$ for every $f\in J$.
Then $C([T^m]f)\in J$ for every $f\in J$ and every $m\geq0$.
\end{lemma}
\begin{proof}
We use strong induction on the degree bound $N=\deg f$, treating the zero
polynomial as having natural degree zero. If $N=0$, then $f=C([T^0]f)$
and all its positive coefficients vanish, giving the assertion.
Suppose $N>0$. The natural degree of $f'$ is strictly less than $N$,
including when $f'=0$. It belongs to $J$ by the stated closure hypothesis.
Induction therefore gives $C([T^{m-1}]f')\in J$ for every $m>0$.
The coefficient derivative formula identifies this polynomial with
$C(m[T^m]f)$. Multiplying it by $C(1/m)$, with the rational inverse
mapped to $R$, proves $C([T^m]f)\in J$ for every $m>0$.
Every positive-degree term $C([T^m]f)T^m$ thus lies in $J$.
The finite sum of all these terms also lies in $J$; subtracting it from
$f\in J$ leaves exactly $C([T^0]f)$. This proves the zero coefficient
as well and completes the induction.
\end{proof}

\begin{theorem}[One-variable differential ideals are coefficient extensions]
\label{hir:differential-stable-polynomial-extension}
\uses{hir:differential-stable-polynomial-coefficients}
\lean{Hironaka.HironakaDifferentialStableCylindricalIdeal.polynomial_eq_map_C_comap_of_derivative_stable}
\leanok
Under the preceding hypotheses, write $C:R\longrightarrow R[T]$ for the
constant-polynomial embedding. Then $J=C_*(C^{-1}J)$.
\end{theorem}
\begin{proof}
The extension of the contraction is contained in $J$, since each of its
generators $C(a)$ belongs to $J$ and $J$ is an ideal. Conversely, let
$f\in J$. The preceding lemma shows that every coefficient of $f$ lies
in $C^{-1}J$. Hence each term $C([T^m]f)T^m$ belongs to its extension.
Their finite sum is $f$, proving the reverse inclusion.
\end{proof}

\begin{theorem}[Differential stability in all finite polynomial variables]
\label{hir:differential-stable-finite-polynomial-extension}
\uses{hir:differential-stable-polynomial-extension}
\lean{Hironaka.HironakaDifferentialStableCylindricalIdeal.mvPolynomial_eq_map_C_comap_of_pderiv_stable}
\leanok
Let $J\subseteq R[T_0,\ldots,T_{d-1}]$ satisfy
$\partial_i f\in J$ for every $i<d$ and every $f\in J$.
For the actual constant-polynomial embedding $C:R\longrightarrow R[T]$,
one has $J=C_*(C^{-1}J)$.
\end{theorem}
\begin{proof}
Induct on $d$. With no variables every polynomial is its constant, and
the assertion follows directly. For $d+1$ variables, single out $T_0$
and regard the same polynomial ring as $S[T_0]$, where
$S=R[T_1,\ldots,T_d]$. The coordinate algebra isomorphism sends
$\partial_0$ to ordinary differentiation in $T_0$: both operations agree
on constants and every variable and obey the additive and product rules.
The preceding theorem therefore identifies $J$ with the extension of its
contraction $J_1$ to $S$.
For every $g\in J_1$ and each remaining variable $T_i$, its constant image
belongs to $J$. Differentiating that image gives the constant image of
$\partial_i g$, which belongs to $J$ by the original hypothesis.
Thus $J_1$ is stable under all its $d$ partial derivatives. Induction
identifies $J_1$ with the extension of its contraction to $R$.
Extensions compose, and the two successive contractions equal contraction
along the composite constant embedding. This yields the claimed equality
for the same original $J$.
\end{proof}

\begin{theorem}[Tangent differential stability gives a genuine normal cylinder]
\label{hir:differential-stable-normal-polynomial-cylinder}
\uses{hir:differential-stable-finite-polynomial-extension}
\lean{Hironaka.HironakaDifferentialStableCylindricalIdeal.ideal_eq_normal_extension_of_tangent_pderiv_stable}
\leanok
Let $J\subseteq\mathbb C[X_0,\ldots,X_{d+t-1}]$. Assume that for every
$i<d$ and \emph{every} $f\in J$, one has $\partial_i f\in J$.
Let $\iota:\mathbb C[Z_0,\ldots,Z_{t-1}]\longrightarrow
\mathbb C[X_0,\ldots,X_{d+t-1}]$ be the genuine embedding
$Z_j\mapsto X_{d+j}$. Then
\[
 J=\iota_*(\iota^{-1}J).
\]
\end{theorem}
\begin{proof}
The coordinate bijection separating the first $d$ and last $t$ variables
identifies the same original polynomial ring with
$\mathbb C[Z_0,\ldots,Z_{t-1}][T_0,\ldots,T_{d-1}]$.
It sends the actual normal embedding $\iota$ to the constant-coefficient
embedding and sends every stated tangent derivative to the corresponding
partial derivative in the $T$ variables. These identities hold on
coefficients and variables, hence on sums and products and therefore on
every polynomial. The image of $J$ consequently satisfies the preceding
theorem over the coefficient $\mathbb Q$-algebra $\mathbb C[Z]$.
Transport its extension--contraction equality back through this genuine
coordinate isomorphism. The resulting contraction and extension are
literally those of $J$ along $Z_j\mapsto X_{d+j}$, proving the formula.
This includes $d=0$ and $t=0$.
\end{proof}

The criterion requires differential stability of the entire ideal. It does
not infer that hypothesis from a chosen family of generators, and it makes
no assertion of analytic normal flatness or of reduced-basis centre powers.
